% TOP_PROBLEM: {"id": 2305020, "title": "Research Problems in Function Theory — Problem 5.20", "category_id": 8, "category": {"id": 8, "name": "analysis", "display_name": "Analysis", "description": "Limits, continuity, calculus, and function theory.", "slug": "analysis", "order_index": 8, "created_at": "2026-07-31T15:26:25.671Z"}, "status": "open"}
% FIRST_POSTED: null
% ENTRY_KIND: statement_only
% Imported from: batch14.tex; UnsolvedMath ID 2305020
% Compile twice: lualatex 2305020.tex
\documentclass[11pt]{article}
\usepackage{fontspec}
\setmainfont{Latin Modern Roman}
\usepackage{amsmath,amssymb,mathtools,unicode-math}
\setmathfont{Latin Modern Math}
\usepackage{xurl,hyperref}
\hypersetup{hidelinks,pdftitle={UnsolvedMath Problem 2305020}}
\newcommand{\sref}[2]{\hyperlink{source:#1:#2}{[S#2]}}
\newcommand{\cataloguescope}[1]{\paragraph{Mathematical question.}#1}
\begin{document}
% UnsolvedMath ID 2305020; source catalogue code AMR-022-5020
\setcounter{section}{2305019}
\section{Research Problems in Function Theory — Problem 5.20}\label{2305020}
{\small\sffamily UnsolvedMath ID 2305020 · Analysis}
\subsection{Definitions and mathematical statement}

\paragraph{Shared notation.}$\mathbb N=\{1,2,\ldots\}$, and $\mathbb Z,\mathbb Q,\mathbb R,\mathbb C$ denote the integers, rationals, reals and complex numbers. Any other conventions are specified locally.


Write $\mathbb D=\{z\in\mathbb C:|z|<1\}$ and $\mathbb T=\partial\mathbb D$. All Taylor coefficients below are taken at $0$. For $\zeta\in\mathbb T$, a Stolz angle means a nondegenerate Euclidean angular sector with vertex $\zeta$, whose two rays enter $\mathbb D$ nontangentially, restricted to a sufficiently small neighborhood of $\zeta$. A finite angular limit is a common finite limit as $z\to\zeta$ within every such sector. Planar measure means two-dimensional Lebesgue measure on the value plane $\mathbb C$; almost every boundary point refers to arc-length measure on $\mathbb T$.
\paragraph{Statement recorded in the supplied source.}
For every holomorphic $f$ on $\mathbb D$, is it true that at almost every $\zeta\in\mathbb T$, either $f$ has a finite angular limit at $\zeta$, or there is a set $E_\zeta\subset\mathbb C$ of planar measure zero such that, for every $a\in\mathbb C\setminus E_\zeta$ and every Stolz angle at $\zeta$, there is a sequence of distinct points $z_j$ in that angle tending to $\zeta$ with $f(z_j)=a$? More generally, what strengthening of a merely dense set of infinitely assumed values is valid in the second alternative? \sref{2305020}{2}
\subsection{Short English statement}
Can the dense set of repeatedly assumed values in Plessner’s boundary alternative be enlarged to all values outside a planar null set?
\subsection{Sources}
\begin{description}
\item[\hypertarget{source:2305020:1}{S1}] ULAM.AI and the UnsolvedMath Contributors, UnsolvedMath: A Curated Collection of Open Mathematics Problems, record 2305020 (AMR-022-5020). CC BY 4.0. \url{https://huggingface.co/datasets/ulamai/UnsolvedMath}.
\item[\hypertarget{source:2305020:2}{S2}] W. K. Hayman and E. F. Lingham, Research Problems in Function Theory (2018), arXiv:1809.07200, Chapter 5, Problem 5.20, its update and the preceding shared notation. \url{https://arxiv.org/abs/1809.07200}.
\end{description}
\label{end:2305020}
\end{document}
